% ======== 编码与语言 ========
% \usepackage[utf8]{inputenc}   % 现代 TeX Live 默认 UTF-8，可省
\usepackage[english]{babel}

% ======== 数学 ========
\usepackage{amsmath, amssymb, amsthm, mathtools, mathrsfs}
\usepackage{dsfont}             % \mathds{R} 等
% 注意：删掉 bbold，与 dsfont 重复

% ======== 图表 ========
\usepackage{graphicx}
\usepackage[export]{adjustbox}
\usepackage{caption, subcaption}
\usepackage{wrapfig}
\usepackage{booktabs}
\usepackage{array, multirow}
\usepackage{tabularx, xltabular}
\usepackage{makecell}
\usepackage{nicematrix}

% ======== TikZ / 绘图 ========
\usepackage{tikz}
\usepackage{pgfplots}
\pgfplotsset{compat=1.18}   % 建议显式设置版本

% ======== 算法 ========
\usepackage{algorithm}
\usepackage{algpseudocodex}

% ======== 列表与排版 ========
\usepackage{microtype}
\usepackage{enumitem}
\usepackage{setspace}
\usepackage{indentfirst}
\usepackage{numprint}
\usepackage{balance}
\usepackage[normalem]{ulem}
\usepackage{soul}
\usepackage{romannum}
\usepackage{xspace}

% ======== 颜色与盒子 ========
\usepackage{xcolor}
\usepackage{tcolorbox}
\tcbuselibrary{skins, breakable, theorems}

% ======== 其它 ========
\usepackage{afterpage}
\usepackage{textcomp}
\usepackage{verbatim}

\usepackage{pifont}
\newcommand{\cmark}{\ding{51}}  % 绿色对勾
\newcommand{\xmark}{\ding{55}}  % 红色叉号

% ======== 自定义 tcolorbox ========
\definecolor{promptDarkGreen}{RGB}{0, 128, 0}
\definecolor{promptLightGreen}{RGB}{240, 250, 240}
\definecolor{inputGray}{gray}{0.30}

\newtcolorbox{examplebox}[1][]{
    enhanced, colback=gray!10, colframe=gray!75,
    sharp corners, boxrule=0.5pt, title=#1
}

\newtcolorbox{promptbox}[1]{
    enhanced, breakable,
    colback=promptLightGreen, colframe=promptDarkGreen,
    coltitle=white, fonttitle=\bfseries\large,
    title={#1},
    sharp corners=south,
    boxrule=1mm, arc=1mm,
    width=\linewidth
}

\newtcolorbox{userinputbox}{
    colback=white, colframe=gray!30,
    boxrule=0.5mm, sharp corners,
    left=2mm, right=2mm, top=2mm, bottom=2mm
}